\providecommand{\keywords}[1]
{
  \small
  \textbf{\textit{Palabras clave---}} #1
}

\phantomsection
% \addcontentsline{toc}{chapter}{Resumen}
\setcounter{page}{\value{abstractpage}}
\begin{center}
  \begin{figure}[h]
    \centering
    \includegraphics[width=0.5\textwidth]{ucab-logo.png} \\
    \centering
    \renewcommand{\baselinestretch}{1.5}
    \membrete
    {
      \textbf{\titulo}
    }
  \end{figure}
  \begin{table}[h!]
    \onehalfspacing
    \raggedleft
    \begin{tabular}{r l}
      Autor:           & \estudianteUno \\
                       & \estudianteDos \\
      Tutor Académico: & \tutor         \\
      Fecha:           & \fecha         \\
    \end{tabular}
  \end{table}
  \vspace{-0.8em}
  \noindent\textbf{Resumen}
  \vspace{-0.4em}
\end{center}

\noindent \small La formación del pensamiento computacional en niños y jóvenes se ha convertido en una necesidad educativa fundamental en la era digital. Sin embargo, las herramientas existentes para esto se basan, mayormente, en el paradigma de flujo de control, lo que genera modelos mentales secuenciales y dificulta la comprensión de conceptos como la concurrencia y los sistemas distribuidos. Ante esta carencia, el presente trabajo tuvo como objetivo desarrollar NodeLand, un lenguaje de programación visual basado en flujos de datos, y su entorno, orientado a fomentar el pensamiento computacional en niños y jóvenes de 8 a 16 años. La propuesta se fundamentó en las definiciones de pensamiento computacional de Wing, Brennan y Resnick, las etapas de desarrollo cognitivo de Piaget, el construccionismo de Papert, y los conceptos de lenguajes basados en dataflow. Metodológicamente, se enmarcó como investigación proyectiva, desarrollada mediante un modelo incremental en cuatro iteraciones. Como resultado, se obtuvo una plataforma web con un editor visual, ejecución concurrente por disponibilidad de datos, un protocolo de distribución hacia tarjetas micro:bit y computadoras, un módulo de desafíos adaptado a dos rangos de edad, y una plataforma social para compartir y remezclar proyectos, y crear extensiones del lenguaje. La validación con expertos confirmó el alto potencial didáctico de la herramienta, y las pruebas de usabilidad demostraron que la representación de programas como grafos de conexiones lógicas resulta intuitiva y asimilable, evidenciando que la dificultad históricamente atribuida a la programación concurrente radica, al menos en parte, en la notación y no en el concepto.

\keywords{pensamiento computacional, programación concurrente, programación visual, dataflow, sistemas distribuidos}
